\section{Deployment}
\label{sec:Deployment}

Durga-generated projects are plain Java applications that connect to Kafka.
They can be deployed as JARs, Docker containers, or Kubernetes workloads
with no Durga-specific runtime dependencies beyond a Kafka cluster.

\subsection{Architecture}

Each generated process instance is a set of Quarkus-based workers (one per
BPMN task) communicating through Kafka topics. There is no central orchestrator
at runtime --- workers consume from input topics and produce to output topics
independently. This means:

\begin{itemize}
\item every worker is horizontally scalable via Kafka consumer groups,
\item deployment is stateless (all state lives in Kafka topics and Streams stores),
\item the generated shaded JAR is self-contained (all dependencies bundled).
\end{itemize}

\subsection{Local development \& testing}

The recommended local loop uses Docker Compose for Kafka and direct JVM
execution for the pipeline:

\begin{verbatim}
# Start Kafka
cd setup && docker compose up -d

# Build and run the pipeline
cd generated/<project>
START_KAFKA=false BOOTSTRAP=localhost:9094 ./run-local.sh
\end{verbatim}

For iterative development, keep Kafka running and rebuild the pipeline:

\begin{verbatim}
mvn -q clean package -DskipTests
java -cp target/<project>-all.jar \
  org.gautelis.durga.generated.<Process>Starter localhost:9094
\end{verbatim}

To test with the monitoring dashboard:

\begin{enumerate}
\item Start Kafka as above.
\item Start the monitoring app:
  \begin{verbatim}
  cd setup
  START_KAFKA=false ./demo-monitoring.sh
  \end{verbatim}
\item Run pipeline scenarios via the generated helper scripts.
\item Open \texttt{http://localhost:18081/dashboard} for the Svelte dashboard
      or \texttt{http://localhost:8081} for the Quarkus REST API.
\end{enumerate}

\subsection{Connection to Kafka}

Kafka bootstrap servers are configured through a chain of fallbacks:

\begin{enumerate}
\item JVM system property \texttt{kafka.bootstrap.servers}.
\item Environment variable \texttt{KAFKA\_BOOTSTRAP\_SERVERS}.
\item Hardcoded default \texttt{localhost:9094}.
\end{enumerate}

In Docker and Kubernetes, set \texttt{KAFKA\_BOOTSTRAP\_SERVERS} as a container
environment variable. The generated \texttt{Dockerfile} passes it to the JVM:

\begin{verbatim}
ENTRYPOINT ["java", "-Xms64m", "-Xmx256m", \
  "-Dkafka.bootstrap.servers=${KAFKA_BOOTSTRAP_SERVERS:-localhost:9094}", \
  "-jar", "/app/pipeline.jar"]
\end{verbatim}

\subsection{Topic provisioning}

Kafka topics must exist before the pipeline starts. There are three approaches,
each supported by the generated artifacts.

\subsubsection{Pre-creation via shell script}

The generated \texttt{topics.sh} creates all required topics using the Kafka
CLI tools bundled with the Kafka distribution:

\begin{verbatim}
KAFKA_BOOTSTRAP_SERVERS=kafka-broker:9092 ./topics.sh
\end{verbatim}

This is suitable for development, staging, and manual production setups.
For automated environments, run \texttt{topics.sh} as an init container or
a CI/CD pipeline step before deploying the workers.

\subsubsection{Init container}

For Kubernetes deployments where the Kafka broker shares the cluster,
add an init container to the generated \texttt{k8s.yml}:

\begin{verbatim}
initContainers:
  - name: topic-init
    image: bitnami/kafka:latest
    command: ['/bin/sh', '-c']
    args:
      - |
        /opt/bitnami/kafka/bin/kafka-topics.sh --create \
          --bootstrap-server $KAFKA_BOOTSTRAP_SERVERS \
          --topic <name> --partitions 3 --replication-factor 1
\end{verbatim}

\subsubsection{Kafka operator (Strimzi)}

In Kubernetes environments running the Strimzi operator, topics can be
declared as \texttt{KafkaTopic} custom resources. A generated
\texttt{kafka-topics.yml} CRD manifest can be applied alongside the
Deployment:

\begin{verbatim}
apiVersion: kafka.strimzi.io/v1beta2
kind: KafkaTopic
metadata:
  name: <process>_<task>_input
  labels:
    strimzi.io/cluster: my-kafka
spec:
  partitions: 3
  replicas: 1
---
\end{verbatim}

The scaffolder currently generates \texttt{topics.sh}; Strimzi CRD
generation can be added as an optional output format.

\subsection{Test environment deployment}

A Docker Compose-based test environment is the fastest way to validate a
pipeline end-to-end. The repository's \texttt{setup/docker-compose.yml}
provides a single-node Kafka broker and Kafka UI.

Add a \texttt{docker-compose.test.yml} to the generated project:

\begin{verbatim}
version: "3"
services:
  pipeline:
    build: .
    environment:
      KAFKA_BOOTSTRAP_SERVERS: kafka:9092
    depends_on:
      - kafka
  kafka:
    image: bitnami/kafka:latest
    environment:
      KAFKA_CFG_NODE_ID: 1
      KAFKA_CFG_PROCESS_ROLES: broker,controller
      KAFKA_CFG_LISTENERS: PLAINTEXT://:9092,CONTROLLER://:9093
      KAFKA_CFG_ADVERTISED_LISTENERS: PLAINTEXT://kafka:9092
      KAFKA_CFG_CONTROLLER_LISTENER_NAMES: CONTROLLER
      KAFKA_CFG_CONTROLLER_QUORUM_VOTERS: 1@kafka:9093
\end{verbatim}

Run with:
\begin{verbatim}
docker compose -f docker-compose.test.yml up --build
\end{verbatim}

\subsection{Production deployment}

The generated \texttt{deploy.sh} supports three targets:

\begin{description}
\item[JAR] \texttt{./deploy.sh} --- builds the shaded JAR and runs locally.
  Suitable for VMs and bare-metal deployments.
\item[Docker] \texttt{DOCKER=true ./deploy.sh} --- builds a Docker image
  tagged \texttt{<processId>:latest}.
\item[Kubernetes] \texttt{DOCKER=true K8S=true ./deploy.sh} --- builds the
  image and applies the generated \texttt{k8s.yml}.
\end{description}

\subsubsection{Docker image}

The generated \texttt{Dockerfile} uses \texttt{eclipse-temurin:21-jre-alpine},
runs as a non-root user, and includes a process-based health check:

\begin{verbatim}
FROM eclipse-temurin:21-jre-alpine
RUN addgroup -S durga && adduser -S durga -G durga
USER durga
COPY target/*-all.jar /app/pipeline.jar
HEALTHCHECK --interval=30s ...
ENTRYPOINT ["java", "-Xms64m", "-Xmx256m", \
  "-Dkafka.bootstrap.servers=${KAFKA_BOOTSTRAP_SERVERS:-...}", \
  "-jar", "/app/pipeline.jar"]
\end{verbatim}

\subsubsection{Kubernetes}

The generated \texttt{k8s.yml} provides:

\begin{itemize}
\item A \texttt{Deployment} with 2 replicas (configurable), resource requests
  and limits (128Mi/512Mi memory, 100m/500m CPU),
\item A \texttt{Service} exposing port 8080,
\item Exec-based liveness and readiness probes.
\end{itemize}

Scaling is done via \texttt{kubectl scale} or the replicas field:

\begin{verbatim}
kubectl scale deployment/<processId> --replicas=4
\end{verbatim}

Maximum parallelism per consumer group is bounded by the number of Kafka
partitions for each topic. Configure partitions in \texttt{topics.sh} before
provisioning, then scale workers up to the partition count.

\subsection{Monitoring in production}

The monitoring topology (\texttt{ProcessMonitoringApp}) is a separate
deployment that consumes the canonical \texttt{process-events} topic and
materializes queryable state. In production:

\begin{itemize}
\item Deploy the monitoring app as a companion service alongside the pipeline
  (same Kafka cluster, different consumer group).
\item The Quarkus REST API and Svelte dashboard provide operational visibility ---
  instance counts, latency summaries, stuck-instance detection, trend reporting.
\item The monitoring topology uses global tables: every instance has a full
  local replica of the monitoring model. Any instance can answer all queries
  without cross-instance RPC.
\end{itemize}

\subsection{Configuration reference}

\begin{longtable}{|>{\raggedright\arraybackslash}p{0.32\textwidth}|>{\raggedright\arraybackslash}p{0.18\textwidth}|>{\raggedright\arraybackslash}p{0.42\textwidth}|}
\hline
\rowcolor[gray]{0.9}\textbf{Variable} & \textbf{Default} & \textbf{Description} \\
\hline
\texttt{KAFKA\_BOOTSTRAP\_SERVERS} & \texttt{localhost:9094} & Kafka broker list \\
\texttt{HTTP\_PORT} & \texttt{8080} & Quarkus HTTP port \\
\texttt{PROFILE} & \texttt{dev} & Quarkus configuration profile \\
\texttt{START\_KAFKA} & \texttt{false} & Start Docker Compose Kafka before running \\
\texttt{DOCKER} & \texttt{false} & Build Docker image in deploy script \\
\texttt{K8S} & \texttt{false} & Apply Kubernetes manifests in deploy script \\
\hline
\end{longtable}
